\startmode[single]

\setvariables[meta]
  [title={Sachtextanalyse},
   subtitle={Die Geschichte vom \quotation{schlechten} Deutsch},
   date={2024-11-27},
   author={Niklas von Hirschfeld},
  ]
\setupinteraction[title={\getvariable{meta}{title}}]

\starttext

\startfrontmatter
\component c_titlepage
\component c_toc
\stopfrontmatter

\setuppagenumber[number=1]

\startbodymatter
\stopmode

\chapter{Notizen}

\startitemize[packed]
\item auch \quotation{grammatisch korrektes} Deutsch kann \quotation{schlecht} sein.
\item die meisten Sprachwissenschaftler lehnen die Differenztheorie ab
\item den richtigen Sprachgebrauch gibt es nicht
\item Verschiedene Arten von Sprachgebrauch. Diese sind entweder:
\startitemize[packed]
\item Funktional angemessen
\item aufgrund sozialer Norm erwartet
\stopitemize
\item viele tendieren dazu Varianten zu bewerten
\startitemize[packed]
\item Dies scheint Sprachen übergreifend
\item Zudem wird diese Berwertung auch auf den Menschen projeziert
\stopitemize
\item In der Bibel:
\startitemize[packed]
\item "falsche" Aussage des Wortes "Schibboleth" identifiziert Mitglied eines feindlichen Stammes
\stopitemize
\item Durch Sprache urteilen wir über:
\startitemize[packed]
\item Schichtzugehörigkeit
\item regionale Herkunft
\item politische Bekenntnisse
\item Charaktereigenschaften
\stopitemize
\item Sprache fungiert als Sozialsymbol
\item subjektive Aspekte und Bewertungen 
\item \quotation{Objektive} Bewertung kann nur innerhalb eines spezifischen soziohistorischen Kontexts untersucht werden
\item \quotation{Standardsprache} ist nicht standartisiert
\stopitemize

\chapter{Sachtextanalyse}

Winfred v. Davies setzt sich in seinem Artikel \quotation{Die Geschichte vom \quotation{schlechten}
Deutsch} kritisch mit der Bewertung und Differenzierung von Sprachvarianten
auseinander. Der Text beleuchtet die Idee, dass die Qualität von Sprache nicht
objektiv festgelegt werden kann, sondern stark von sozialen und funktionalen
Kontexten abhängt. Im Folgenden wird die Methodik der Argumentation analysiert,
um die zentralen Aussagen des Autors sowie deren sprachwissenschaftliche und
gesellschaftliche Relevanz herauszuarbeiten.

Davies beginnt seinen Text mit der Feststellung, dass auch grammatisch
korrektes Deutsch als „schlecht“ empfunden werden kann. Damit stellt er die
gängige Annahme in Frage, dass Sprachqualität allein an formalen grammatischen
Regeln gemessen werden kann. Diese These dient als Grundlage für die weitere
Diskussion und wird durch zahlreiche Beispiele und sprachwissenschaftliche
Ansätze untermauert.

Anschließend erläutert er, dass die meisten sprachwissenschaftler die
\quotation{Differenztheorie} ablehnen. Diese besagt, dass alle Variationen in
ihrerer Ausdrucksmöglichkeit und logischen Analysekapazität einer funktionalen
äquivalent sind.

\startmode[single]

\stopbodymatter

\startbackmatter
\component c_definitions
% \component c_bibliography
\stopbackmatter

\stoptext
\stopmode
